% TOP_PROBLEM: {"problemId": "problem.bass-quillen-conjecture", "canonicalTitle": "Bass-Quillen conjecture", "releaseRank": 312, "primaryDomain": "algebra_representation_category", "primaryDomainLabel": "Algebra, representation & category theory", "displayStatus": "open_with_solved_subcases", "releaseStatus": "open_with_solved_subcases"}
% !TeX program = lualatex
% ATTEMPT_STATUS: unresolved
\documentclass[11pt]{article}
\usepackage[margin=25mm]{geometry}
\usepackage{fontspec}
\setmainfont{Latin Modern Roman}
\usepackage{amsmath,amssymb,mathtools}
\usepackage{unicode-math}
\setmathfont{Latin Modern Math}
\usepackage{xurl,hyperref}
\hypersetup{colorlinks=true,urlcolor=blue}
\setcounter{secnumdepth}{0}
\setlength{\parindent}{0pt}
\setlength{\parskip}{5pt}
\setlength{\emergencystretch}{3em}
\begin{document}
\section*{\texorpdfstring{Bass-Quillen conjecture}{Bass-Quillen conjecture}}\label{312}
\subsection{Definitions and mathematical statement}
Let $R$ be a commutative regular Noetherian ring: each localization $R_{\mathfrak p}$ is a regular local ring, whose maximal ideal needs $\dim R_{\mathfrak p}$ generators. A module is finitely generated projective if it is a direct summand of a finite free module. The Bass--Quillen assertion says that for every $d\ge0$ and finitely generated projective $R[t_1,\ldots,t_d]$-module $P$, there exists a finitely generated projective $R$-module $P_0$ with
\[
 P\cong P_0\otimes_R R[t_1,\ldots,t_d].
\]
One may take $P_0=P/(t_1,\ldots,t_d)P$. In geometric language, every vector bundle on affine $d$-space over $R$ is pulled back from $\operatorname{Spec}R$. The Noetherian regular-ring convention is that of the cited formulation; a broader meaning of ``regular'' needs separate specification.
\par\smallskip\textit{Formulation and concept references:} \href{https://arxiv.org/html/2201.06424v4}{[S1]}.

\subsection{Short English statement}
Over a regular ring, must every finite-rank vector bundle on affine space come from a vector bundle on the base ring?
\subsection{Research attempt: polynomial descent versus formal triviality}
\paragraph{1. Scope and source audit.}
The conclusion is $P\cong P_0\otimes_R R[t_1,\ldots,t_d]$ with $P_0=P/(t_1,\ldots,t_d)P$. It is not freeness over an arbitrary regular base ring. \v{C}esnavi\v{c}ius's survey [S1], Section 2.1, was read on 27 September 2026. It gives the local and one-variable reductions, records known unramified and line-bundle cases, and identifies ramified regular local rings as the main outstanding setting. It also explicitly warns that a purported general deduction for excellent Henselian local rings in older references is incorrect. A 2025 paper of Guo and Liu [E1] treats torsors over smooth algebras over valuation rings with additional group hypotheses; that scope is not every ramified regular ring. The proof below isolates why two tempting elementary shortcuts do not settle the general conjecture.

\paragraph{2. The special fiber is projective and descent is uniquely identified.}
Choose a complement $Q$ such that $P\oplus Q\cong R[\mathbf t]^m$. Tensor this identity with $R$ via $t_i\mapsto0$. It gives
\[
 P_0\oplus Q_0\cong R^m,
\]
so $P_0$ is finite projective. If $P\cong A\otimes_RR[\mathbf t]$ for some $R$-module $A$, then specializing at zero gives $A\cong P_0$. Thus the question is existence of a descent isomorphism; the isomorphism class of a possible descended module is already determined. It would not suffice to construct a projective module of the same rank or the same $K_0$ class.

If $R$ is local, every finite projective $R$-module is free. Here is the standard short proof. Lift a basis of $P_0/\mathfrak mP_0$ to generators of $P_0$ using Nakayama. This gives a surjection $R^r\to P_0$, split because $P_0$ is projective. Its finite kernel $K$ satisfies $K/\mathfrak mK=0$, because reduction of the splitting has dimensions $r=r+\dim(K/\mathfrak mK)$. Nakayama gives $K=0$. Hence the conjecture over a local regular ring asks that every finite projective module over its polynomial ring be free. The polynomial ring itself is not local; the same one-line argument cannot be applied directly to it.

\paragraph{3. One-variable reduction and an elementary base case.}
Suppose the one-variable assertion has been proved for every regular Noetherian ring. Polynomial rings over such rings remain regular Noetherian. Given $P$ over $R[t_1,\ldots,t_d]$, apply the assertion with base $R[t_1,\ldots,t_{d-1}]$ and variable $t_d$. This descends $P$ to its $t_d=0$ fiber. Repeat, finally obtaining $P_0$. Thus one variable is enough, provided the theorem is known uniformly for every regular base ring encountered. Proving one variable only for fields does not permit this induction, because the new base rings are no longer fields.

For a field $k$ and one variable, $k[t]$ is a principal ideal domain. A finite projective module is a direct summand of a finite free module, hence a submodule of a finite free module. The usual Euclidean-algorithm induction proves every such submodule free: project to one coordinate, choose a generator of the image ideal, split off one lift, and apply induction to the kernel in one fewer coordinate. Therefore $P$ is free, with rank equal to that of its specialization. This proves the one-variable field subcase. The all-variable field theorem is Quillen--Suslin, a deeper established input, not supplied by repeating the PID argument in several variables.

The local-to-global reduction in [S1] uses Quillen patching. Local descent isomorphisms need not agree on overlaps automatically; treating their existence as if they were already a compatible cocycle omits the patching theorem.

\paragraph{4. Formal descent holds without regularity.}
Represent a projective module over $R[t]$ as the image of an idempotent matrix $E(t)\in M_m(R[t])$, and put $E_0=E(0)$. Define
\[
 U(t)=E(t)E_0+(I-E(t))(I-E_0).
\]
The idempotent identities imply
\[
 U(t)E_0=E(t)E_0=E(t)U(t),\qquad U(0)=I.
\]
Over $R[[t]]$, any matrix with constant term $I$ is invertible: if $U=I+V$ with $V\in tM_m(R[[t]])$, the formal geometric series $\sum_{j\geq0}(-V)^j$ is a well-defined inverse because each coefficient receives only finitely many contributions. Hence
\[
 E(t)=U(t)E_0U(t)^{-1}\quad\text{over }R[[t]].
\]
It follows that every polynomial projective module becomes the extension of its zero fiber after this formal base change. The argument works for every commutative ring $R$, with no regularity assumption. It cannot therefore be the missing proof of the regular-ring conjecture: it establishes a weaker formal assertion.

The same formula works in several variables after completion in the ideal $(t_1,\ldots,t_d)$. It also gives exact conjugacy modulo every power of that ideal. Having compatible solutions to all finite-order problems does not show that one solution and its inverse are polynomial.

\paragraph{5. An explicit idempotent demonstrates the algebraization gap.}
Over $\mathbb Z[t]$, set
\[
 E(t)=\begin{pmatrix}1-t^2&t\\t-t^3&t^2\end{pmatrix},
 \qquad E_0=\begin{pmatrix}1&0\\0&0\end{pmatrix}.
\]
Write $E=vw$ with $v=(1,t)^T$ and $w=(1-t^2,t)$. Since $wv=1$, $E^2=E$. The canonical matrix from Section 4 is
\[
 U=\begin{pmatrix}1-t^2&-t\\t-t^3&1-t^2\end{pmatrix},
 \qquad\det U=1-t^2.
\]
Its determinant is a unit in $\mathbb Z[[t]]$ but not in $\mathbb Z[t]$. Therefore $U$ is not an invertible polynomial matrix even though its value at zero is the identity and it solves the intertwining identity.

This example is not a counterexample to Bass--Quillen. In fact
\[
 V=\begin{pmatrix}1&-t\\t&1-t^2\end{pmatrix},\qquad\det V=1,
 \qquad E=VE_0V^{-1}.
\]
The first column of $V$ is $v$, and the second lies in the kernel of $w$, which verifies the conjugacy. The example sharply separates two statements: a polynomial conjugator may exist, while the obvious formal conjugator need not be polynomially invertible. A general proof must construct or descend an appropriate conjugator, rather than assume the first one works.

\paragraph{6. Stable information cannot be silently cancelled.}
A unimodular row $a=(a_1,\ldots,a_m)$ over a ring $A$ defines a surjection $A^m\to A$. Choosing a column $b$ with $ab=1$ splits it, so its kernel $K$ satisfies
\[
 K\oplus A\cong A^m.
\]
The row can be completed to an invertible matrix precisely when this kernel is free of rank $m-1$: a basis of the kernel together with $b$ gives such a matrix, and conversely the remaining columns of an inverse completion provide a basis. Thus cancellation of the extra free summand is a real mathematical issue.

For a concrete regular base where it fails, let
\[
 R=\mathbb R[x,y,z]/(x^2+y^2+z^2-1),
 \qquad K=\ker\bigl(R^3\xrightarrow{(x,y,z)}R\bigr).
\]
The row is unimodular because $x^2+y^2+z^2=1$, so $K\oplus R\cong R^3$. The affine quadric is smooth, since its gradient cannot vanish on the defining hypersurface, and its coordinate ring is regular. On real points, the associated vector bundle is the tangent bundle of $S^2$. If $K$ were free, evaluating an algebraic basis would trivialize that real tangent bundle, contradicting the hairy-ball theorem. Thus $K$ is stably free but not free.

This is an example over the base ring, not a counterexample to polynomial descent. The extended module $K\otimes_RR[t]$ is exactly the type of module that the conjecture permits. The example explains why global freeness is the wrong target and why equality of stable $K$-theory classes alone does not establish an actual module isomorphism.

\paragraph{7. Rank one and the mixed-characteristic obstruction.}
The line-bundle case is already known: [S1], Section 2.1.2(8), records homotopy invariance of the Picard group for seminormal rings, hence for regular rings. In the local regular case it can also be understood through factoriality: a regular local ring is a UFD, its polynomial ring remains a UFD, and invertible rank-one modules have trivial divisor class. These are established commutative-algebra theorems. They settle rank one, but taking determinants of a higher-rank module only gives its line-bundle invariant and cannot reconstruct the module.

For ramified mixed-characteristic regular local rings, formal conjugacy, determinant triviality, and stable homotopy invariance all leave the same kind of descent or cancellation issue. The unramified methods described in [S1] have structural hypotheses that cannot simply be asserted after completion. In particular, completeness does not turn a ramified ring into an unramified one.

\paragraph{8. Diagnostics and conclusion.}
The saved exact polynomial-matrix calculation verifies $E^2=E$, $UE_0=EU$, both determinants, and $VE_0V^{-1}=E$. It also checks the unimodular-row projector at 21 exact rational points of the sphere; this finite check is not a symbolic quotient-ring proof. These checks concern the displayed examples; they neither classify all idempotent matrices nor certify algebraization of arbitrary formal solutions.

\textbf{UNRESOLVED.} The attempt proves elementary field and formal subcases, the one-variable reduction, and explicit diagnostics separating formal, stable, and actual polynomial descent. The first general gap is an algebraic descent/cancellation argument for arbitrary higher-rank projectives over the polynomial ring of a ramified regular local base, followed by the established patching framework.

\subsection{Sources}
\begin{itemize}
\item[S1]Kestutis Cesnavicius, Problems about torsors over regular rings with an appendix by Yifei Zhao, arXiv:2201.06424v4, 2025.
\url{https://arxiv.org/html/2201.06424v4}.
\textit{Formulation source.}
\item[E1]N. Guo and F. Liu, The Bass--Quillen conjecture for torsors over valuation rings, arXiv:2503.10163.
\url{https://arxiv.org/abs/2503.10163}.
\textit{Primary abstract checked 27 September 2026. Smooth algebras over valuation rings and specified reductive groups are the stated scope, not all ramified regular local rings.}
\end{itemize}
\end{document}
